\scanpage{27}
\uline{Corollaire}\qquad $G^+\xrightarrow{\sim}\mathrm{Alt}_T$

Reconstruction de l'icosaèdre \uline{orienté} par l'ens orienté $T$

Soit $\Pi$ icosaèdre, \deleted{\illegible{1}} pas orienté pour le moment. L'ens $T$ de ses trièdres distingués est orienté de façon canonique, soit $\unsure{\bar\omega|\varpi}$ son orientation. Sur $T\deleted{\unsure{\{}}$ il y a exactem$\unsure{^{\mathrm t}|t}$ $\frac12\,4!$ \inserted{$=12$} structures polygonales, qui se partagent en $2$ paquets de \inserted{$6$} structures polygonales associées à l'une et l'autre orientation. Soit $\mathrm{Pent}^+(T,\bar\omega)$ l'ens des \deleted{\unsure{$1$}}$6$ structures pentagonales associées à $\bar\omega$. On a une application canonique
\[
S/\deleted{\unsure{G}}\underline a\longrightarrow\mathrm{Pent}^+(T,\bar\omega)
\]
compatible avec op. de $G$. Or $G$ opère transitivement sur $S$, et il est immédiat que $\mathrm{Alt}_T$ opère transitivemt sur $\mathrm{Pent}^+(T,\bar\omega)$ (valable pour tt \unsure{$a$-|$n$-}polygône à $n\equiv1\ (4)$ côtés \ldots), donc l'application précédente est surjective, donc bijective
\[
S/\underline a\simeq\mathrm{Pent}^+(T,\bar\omega).
\]
% archive sheet number: 26
